\documentclass[11pt]{article}

\usepackage[T1]{fontenc}
\usepackage{lmodern}
\usepackage[margin=1.08in]{geometry}
\usepackage{amsmath,amssymb,amsthm,mathtools}
\usepackage{booktabs}
\usepackage{enumitem}
\usepackage{xcolor}
\usepackage{microtype}
\usepackage[hidelinks]{hyperref}

\definecolor{wallred}{RGB}{142,31,31}
\newtheorem{theorem}{Theorem}[section]
\newtheorem{proposition}[theorem]{Proposition}
\newtheorem{corollary}[theorem]{Corollary}
\newtheorem{lemma}[theorem]{Lemma}
\newtheorem{remark}[theorem]{Remark}
\newtheorem{example}[theorem]{Example}
\newtheorem{definition}[theorem]{Definition}
\newcommand{\TV}{\mathrm{TV}}
\newcommand{\osc}{\operatorname{osc}}
\newcommand{\diamH}{\Delta_{H}}
\newcommand{\R}{\mathbb{R}}
\newcommand{\1}{\mathbf{1}}

\title{\textbf{Local Tanh, Global Wall}\\[2mm]
Tensorization, Equality, and the Birkhoff--Dobrushin Divide}
\author{The Eriksson Programme}
\date{August 4, 2026}

\begin{document}
\maketitle

\begin{abstract}
For a strictly positive finite Markov kernel $P$, Birkhoff's projective
contraction coefficient is $\tanh(\diamH(P)/4)$, whereas Dobrushin's
coefficient is the exact contraction coefficient for oscillation.  The sharp
comparison $\delta(P)\leq \tanh(\diamH(P)/4)$ is known.  We give a finite,
explicit equality classification: apart from the rank-one case, equality
holds exactly when a diameter-realizing pair of rows has a likelihood ratio
taking two reciprocal values.  We then record the exact tensor law
$\diamH(P\otimes Q)=\diamH(P)+\diamH(Q)$.  Consequently the global Birkhoff
coefficient of $P^{\otimes L}$ obeys hyperbolic addition and tends to one for
every nonzero single-site diameter.  The global Dobrushin coefficient also
tends to one in every nontrivial product, while each coordinate oscillation
contracts with the exact, volume-independent factor $\delta(P)$.  For the
binary Ising link all local quantities coincide at $\tanh |J|$, but the
global projective coefficient is $\tanh(L|J|)$.  Thus the same calculation
both explains the ubiquitous local $\tanh$ and supplies a sharp wall against
deducing a volume-uniform gap from the global tensor cone.  We separate
classical inputs, elementary tensor consequences, and the narrow equality
and localization synthesis, and we include a reproducible exact-arithmetic
kill test.  A companion Lean module machine-checks the algebraic core of the
equality, tensorization, hyperbolic-addition, and coordinate-averaging steps;
the classical analytic Birkhoff--Hopf input remains external.
\end{abstract}

\section{Introduction}

Two contraction mechanisms are often placed side by side in finite-volume
statistical mechanics.  Birkhoff's theorem measures contraction in Hilbert's
projective metric on a positive cone.  Dobrushin's coefficient measures
contraction of oscillations, or equivalently total variation between rows of
a Markov kernel.  Both produce a hyperbolic tangent in the binary Ising
cell.  It is tempting to read this coincidence as evidence that the global
projective calculation supplies a volume-uniform mixing constant.  The
temptation is wrong in the most elementary independent product.

The purpose of this paper is to make the coincidence and the obstruction part
of one theorem.  There are three points.

\begin{enumerate}[label=(\roman*)]
\item The sharp total-variation--Hilbert comparison is a pairwise statement.
We record its complete equality condition and lift it to finite kernels.
\item Hilbert diameter is exactly additive under tensor products.  Hence the
Birkhoff coefficient combines by hyperbolic addition, not by taking a maximum.
\item Dobrushin's global coefficient also degenerates on independent products.
The volume-uniform statement lives instead in the vector of coordinate
oscillations and, for interacting specifications, inside the Dobrushin
interdependence window.
\end{enumerate}

The wall is therefore not a caveat after the theorem.  It is one of the
theorem's conclusions.

\subsection*{Claim ledger}

The classical projective contraction formula is due to Birkhoff
\cite{Birkhoff1957} and was developed further by Bushell
\cite{BushellRatio,BushellMetric}.  The Markov coefficient and the
interdependence method originate with Dobrushin
\cite{Dobrushin1956,Dobrushin1968,Dobrushin1970}.  Gaubert and Qu identify
Dobrushin contraction with contraction in Hopf's oscillation seminorm on
cones \cite{GaubertQu2014}.  Reeb, Kastoryano and Wolf prove a general
base-norm versus projective-diameter estimate
\cite{ReebKastoryanoWolf2011}.  Cohen and Fausti give the sharp
total-variation--Hilbert inequality for probability measures, including
sharpness \cite{CohenFausti2024}.  We do not claim any of these statements as
new.

The narrow contribution of the present note is the following package: an
if-and-only-if equality classification for finite positive distributions and
kernels; the exact tensor hyperbolic law placed next to that classification;
and an exact global/local product dichotomy with the binary Ising saturation
as a common test cell.  The tensor identity itself is elementary.  Our
priority search found no source stating this entire finite-kernel package, but
that search is not a claim of terminal priority.

\section{Definitions and conventions}

Let $E$ be a nonempty finite set.  For $u,v\in\R^E_{>0}$, Hilbert's
projective distance is
\begin{equation}\label{eq:hilbert}
 d_H(u,v)=\log\frac{\max_{i\in E}u_i/v_i}
                        {\min_{i\in E}u_i/v_i}.
\end{equation}
It is a metric on rays, and it vanishes exactly when $u$ and $v$ are
proportional.

For a strictly positive matrix $A=(a_{xy})_{x\in X,y\in Y}$, define
\begin{align}
 \Theta(A)&=\max_{\substack{x,x'\in X\\ y,y'\in Y}}
 \frac{a_{xy}a_{x'y'}}{a_{xy'}a_{x'y}},\label{eq:theta}\\
 \diamH(A)&=\log\Theta(A).
\end{align}
Equivalently, $\diamH(A)$ is the largest Hilbert distance between two rows
of $A$.  Birkhoff's formula gives the projective contraction coefficient
\begin{equation}\label{eq:birkhoff}
 \tau(A)=\tanh\!\left(\frac{\diamH(A)}4\right).
\end{equation}
Transposition changes neither $\Theta$ nor $\diamH$, so our row convention
does not affect this number.

A Markov kernel $P$ on $E$ is row stochastic.  Its Dobrushin coefficient is
\begin{equation}\label{eq:dobrushin}
 \delta(P)=\frac12\max_{x,x'\in E}
             \sum_{y\in E}|P(x,y)-P(x',y)|.
\end{equation}
We use $\TV(p,q)=\frac12\sum_y|p_y-q_y|$.  Dobrushin's 1956 overlap
convention is $1-\delta(P)$, so formulas must be translated before comparing
the two notations.  For $f:E\to\R$, let
\[
 \osc(f)=\max_E f-\min_E f.
\]
The standard dual formula says
\begin{equation}\label{eq:osc-exact}
 \delta(P)=\sup_{\osc(f)>0}\frac{\osc(Pf)}{\osc(f)}.
\end{equation}

\section{The sharp pairwise comparison and all equality cases}

The inequality in this section is known in a more general measure-theoretic
form \cite{CohenFausti2024}.  The finite proof is included because it exposes
the equality condition needed later.

\begin{theorem}[Sharp pairwise comparison]\label{thm:pairwise}
Let $p,q$ be strictly positive probability vectors on a finite set.  Put
\[
 r_i=\frac{p_i}{q_i},\qquad m=\min_i r_i,\qquad M=\max_i r_i.
\]
Then
\begin{equation}\label{eq:two-step}
 \TV(p,q)
 \leq \frac{(M-1)(1-m)}{M-m}
 \leq \frac{\sqrt M-\sqrt m}{\sqrt M+\sqrt m}
 =\tanh\!\left(\frac{d_H(p,q)}4\right),
\end{equation}
with the middle expression interpreted as zero when $p=q$.

If $p\ne q$, equality between the first and last terms holds if and only if
there is a number $t>1$ such that
\begin{equation}\label{eq:reciprocal-ratio}
 \left\{\frac{p_i}{q_i}:i\in E\right\}=\{t,t^{-1}\}.
\end{equation}
In that case, with $A=\{i:p_i=tq_i\}$,
\begin{equation}\label{eq:mass-swap}
 q(A)=\frac1{t+1},\quad p(A)=\frac{t}{t+1},\quad
 q(A^c)=\frac{t}{t+1},\quad p(A^c)=\frac1{t+1}.
\end{equation}
Thus the two measures exchange the masses of the two blocks.
\end{theorem}

\begin{proof}
Since $\sum_iq_ir_i=1$, either $p=q$ or $m<1<M$.  For every
$r\in[m,M]$,
\[
 (r-1)_+\leq \frac{M-1}{M-m}(r-m).
\]
Taking expectation under $q$ gives the first inequality because
$\TV(p,q)=\sum_iq_i(r_i-1)_+$.

Write $a=\sqrt M$ and $b=\sqrt m$.  With the positive common denominator
$(a-b)(a+b)$, the numerator of the difference between the second and first
bounds in \eqref{eq:two-step} is
\begin{align*}
 (a-b)^2-(a^2-1)(1-b^2)=(ab-1)^2.
\end{align*}
This proves the second inequality.  The last identity follows from
$d_H(p,q)=\log(M/m)$.

In the nontrivial case, equality in the chord bound requires every $r_i$ to
be an endpoint, so $r_i\in\{m,M\}$.  Equality in the algebraic step requires
$mM=1$.  Hence $M=t$ and $m=t^{-1}$ for $t>1$.  Conversely these two
conditions make both inequalities equalities.  Finally normalization gives
$tq(A)+t^{-1}q(A^c)=1$ together with $q(A)+q(A^c)=1$, which yields
\eqref{eq:mass-swap}.
\end{proof}

\begin{remark}
The equality condition is rigid in likelihood-ratio space but allows an
arbitrary number of atoms in each of the two blocks.  For two-point
probabilities it simply says that the two vectors are nontrivial swaps.
\end{remark}

\section{Kernel comparison and equality classification}

\begin{theorem}[Birkhoff--Dobrushin comparison with equality]
\label{thm:kernel-equality}
Let $P$ be a strictly positive finite Markov kernel.  Then
\begin{equation}\label{eq:kernel-bound}
 \delta(P)\leq \tau(P)=
 \tanh\!\left(\frac{\diamH(P)}4\right).
\end{equation}
If $\diamH(P)=0$, all rows of $P$ agree and both sides vanish.  If
$\diamH(P)>0$, equality holds in \eqref{eq:kernel-bound} if and only if there
are rows $x,x'$ and a number $t>1$ such that
\begin{enumerate}[label=(\alph*)]
\item $d_H(P(x,\cdot),P(x',\cdot))=\diamH(P)$; and
\item $P(x,y)/P(x',y)\in\{t,t^{-1}\}$ for every $y$.
\end{enumerate}
Necessarily $t=e^{\diamH(P)/2}$.
\end{theorem}

\begin{proof}
Apply Theorem~\ref{thm:pairwise} to every pair of rows:
\[
 \TV(P(x,\cdot),P(x',\cdot))
 \leq \tanh\!\left(\frac{d_H(P(x,\cdot),P(x',\cdot))}{4}\right)
 \leq \tau(P).
\]
Taking the maximum proves \eqref{eq:kernel-bound}.  If equality holds, choose
a row pair attaining $\delta(P)$.  Both displayed inequalities must then be
equalities.  The second makes the pair diameter-realizing, and
Theorem~\ref{thm:pairwise} gives the reciprocal two-valued likelihood ratio.
The converse is immediate.  For such a pair, its maximum-to-minimum ratio is
$t^2$, so $2\log t=\diamH(P)$.
\end{proof}

\begin{corollary}[Binary kernels]
For a strictly positive two-state kernel with distinct rows,
$\delta(P)=\tau(P)$ if and only if its rows are $(a,1-a)$ and $(1-a,a)$ for
some $a\ne\frac12$.
\end{corollary}

\begin{proof}
Theorem~\ref{thm:kernel-equality} says the two likelihood ratios are
reciprocal.  Equivalently $a(1-a)=b(1-b)$ for rows $(a,1-a)$ and $(b,1-b)$.
Distinctness forces $b=1-a$.
\end{proof}

The comparison can be strict even in dimension two.  For example, the rows
$(0.9,0.1)$ and $(0.6,0.4)$ have $\delta=0.3$, while their likelihood ratios
$3/2$ and $1/4$ are not reciprocal, so $\delta<\tau$.

\section{The tensor theorem and the projective wall}

\begin{theorem}[Exact tensor law]\label{thm:tensor}
For strictly positive finite matrices $A$ and $B$,
\begin{equation}\label{eq:tensor-diameter}
 \Theta(A\otimes B)=\Theta(A)\Theta(B),\qquad
 \diamH(A\otimes B)=\diamH(A)+\diamH(B).
\end{equation}
Consequently
\begin{equation}\label{eq:hyperbolic-addition}
 \tau(A\otimes B)=
 \frac{\tau(A)+\tau(B)}{1+\tau(A)\tau(B)}.
\end{equation}
In particular,
\begin{equation}\label{eq:tensor-power}
 \diamH(A^{\otimes L})=L\diamH(A),\qquad
 \tau(A^{\otimes L})=
 \tanh\!\left(L\operatorname{arctanh}\tau(A)\right).
\end{equation}
\end{theorem}

\begin{proof}
A cross ratio of $A\otimes B$ is the product of one cross ratio of $A$ and
one cross ratio of $B$.  Hence it is at most $\Theta(A)\Theta(B)$.  Because
all index sets are finite, maximizing indices can be selected independently
in the two factors, so the upper bound is attained.  This proves
\eqref{eq:tensor-diameter}.  Equations \eqref{eq:hyperbolic-addition} and
\eqref{eq:tensor-power} follow from Birkhoff's formula and the addition law
for $\tanh$.
\end{proof}

\begin{corollary}[The global projective wall]\label{cor:wall}
If $\diamH(A)>0$, then $\tau(A^{\otimes L})\uparrow1$.  Therefore direct
application of Birkhoff's theorem to the full positive tensor cone cannot
produce a contraction factor $q<1$ independent of $L$.  If $\diamH(A)=0$,
then $A$ has projective rank one and $\tau(A^{\otimes L})=0$.
\end{corollary}

The statement is deliberately method-specific.  It rules out a
volume-uniform constant obtained from the global tensor-cone diameter.  It
does not rule out a different cone, a quotient adapted to locality, a
weighted nonproduct seminorm, a block dynamics argument, or an independent
spectral method.  Each such change is a genuinely different mechanism.

\begin{proposition}[Diagonal-weight blindness]\label{prop:diagonal}
If $D_X,D_Y$ are positive diagonal matrices, then
\[
 \Theta(D_XAD_Y)=\Theta(A),\qquad \diamH(D_XAD_Y)=\diamH(A).
\]
Thus positive source and target weights cannot repair the tensor wall while
the same projective cross ratios are used.
\end{proposition}

\begin{proof}
Every row factor from $D_X$ and every column factor from $D_Y$ occurs once in
the numerator and once in the denominator of each cross ratio, and cancels.
\end{proof}

\section{Global Dobrushin degeneration and exact localization}

The global wall is not unique to Hilbert's metric.  Independent products
separate distinct product measures, so their global total-variation
coefficient also tends to one.

\begin{theorem}[Global degeneration]\label{thm:global-delta}
Let $P$ be a finite Markov kernel with two distinct positive rows $p$ and
$q$.  Set their Hellinger affinity
\[
 \rho(p,q)=\sum_y\sqrt{p_yq_y}<1.
\]
Then
\begin{equation}\label{eq:global-delta-lower}
 \delta(P^{\otimes L})
 \geq \TV(p^{\otimes L},q^{\otimes L})
 \geq 1-\rho(p,q)^L.
\end{equation}
Consequently $\delta(P^{\otimes L})\to1$.
\end{theorem}

\begin{proof}
The constant input words corresponding to the two rows select
$p^{\otimes L}$ and $q^{\otimes L}$ as rows of the product kernel.  For any
probability vectors $\mu,\nu$,
\[
 1-\TV(\mu,\nu)=\sum_z\min(\mu_z,\nu_z)
 \leq\sum_z\sqrt{\mu_z\nu_z}.
\]
The affinity tensorizes, giving \eqref{eq:global-delta-lower}.  Strict
Cauchy--Schwarz gives $\rho(p,q)<1$ when $p\ne q$.
\end{proof}

The volume-uniform statement is coordinatewise.  On $E^L$ define
\[
 \osc_i(f)=\max_{x_{-i}}\left(
       \max_{a\in E}f(x_{-i},a)-\min_{a\in E}f(x_{-i},a)
                         \right).
\]

\begin{theorem}[Exact local product contraction]\label{thm:local}
For every $f:E^L\to\R$ and every coordinate $i$,
\begin{equation}\label{eq:local-contraction}
 \osc_i(P^{\otimes L}f)\leq\delta(P)\osc_i(f).
\end{equation}
The constant is exact whenever $\delta(P)>0$.  Hence both
$\max_i\osc_i$ and $\sum_i\osc_i$ contract by the same factor $\delta(P)$,
independently of $L$.
\end{theorem}

\begin{proof}
Compare two inputs that differ only at coordinate $i$ and condition on all
output coordinates other than $i$.  For each fixed outside output $z_{-i}$,
the function $a\mapsto f(z_{-i},a)$ has oscillation at most $\osc_i(f)$.
The one-coordinate dual formula \eqref{eq:osc-exact} bounds the resulting
difference by $\delta(P)\osc_i(f)$.  Averaging over $z_{-i}$ and maximizing
over the two inputs proves \eqref{eq:local-contraction}.

For exactness, choose a one-site function $g$ attaining
\eqref{eq:osc-exact} and put $f(x)=g(x_i)$.  Then both sides reduce to the
one-site equality.  The final two assertions follow by taking a maximum or a
sum over $i$.
\end{proof}

There is no contradiction between Theorems~\ref{thm:global-delta} and
\ref{thm:local}: the seminorms are different.  The global oscillation can
compare two words differing at every coordinate; $\osc_i$ permits one
coordinate change at a time and retains the geometry needed by local
comparison theory.

\section{The binary Ising cell: exact coincidence and saturation}

For $J\in\R$, consider
\begin{equation}\label{eq:ising-link}
 K_J=\begin{pmatrix}e^J&e^{-J}\\ e^{-J}&e^J\end{pmatrix},
 \qquad P_J=\frac1{2\cosh J}K_J.
\end{equation}

\begin{proposition}[The common local $\tanh$]\label{prop:ising}
For the binary link,
\begin{equation}\label{eq:ising-coincidence}
 \diamH(K_J)=4|J|,\qquad
 \tau(K_J)=\delta(P_J)=\tanh|J|.
\end{equation}
The nontrivial normalized eigenvalue of $K_J$ is $\tanh J$.  Moreover, for
the binary conditional
\[
 \pi_h(+)=\frac{e^h}{e^h+e^{-h}}=\frac{1+\tanh h}{2},
\]
one has, for $J\geq0$,
\begin{equation}\label{eq:conditional-envelope}
 \sup_{h\in\R}\TV(\pi_{h+J},\pi_{h-J})=\tanh J,
\end{equation}
with equality at $h=0$.
\end{proposition}

\begin{proof}
Diagonal entries against antidiagonal entries give the maximal cross ratio
$e^{4|J|}$.  The two rows of $P_J$ are swaps, and their total variation is
$|e^J-e^{-J}|/(e^J+e^{-J})=\tanh|J|$.  The eigenvectors $(1,1)$ and $(1,-1)$
give normalized eigenvalues $1$ and $\tanh J$.

For the conditional calculation,
\[
 \TV(\pi_{h+J},\pi_{h-J})
 =\frac12|\tanh(h+J)-\tanh(h-J)|.
\]
The difference is even in $h$ and decreases for $h>0$; at $h=0$ it equals
$\tanh J$.
\end{proof}

\begin{corollary}[Ising tensor saturation and wall]
For $J\ne0$,
\begin{equation}\label{eq:ising-wall}
 \tau(K_J^{\otimes L})=\tanh(L|J|)\longrightarrow1,
\end{equation}
while the largest modulus below the top normalized eigenvalue of the product
is $\tanh|J|$ for every $L\geq1$.  The coordinatewise Dobrushin constant is
also $\tanh|J|$ for every $L$.
\end{corollary}

For completeness, put $a=(1+\tanh J)/2$.  Two antipodal rows of
$P_J^{\otimes L}$ induce the laws $\operatorname{Bin}(L,a)$ and
$\operatorname{Bin}(L,1-a)$ after counting plus spins.  Thus their total
variation gives an exact finite sum lower bound for the global coefficient,
and Theorem~\ref{thm:global-delta} yields the transparent estimate
\begin{equation}\label{eq:ising-hellinger}
 \delta(P_J^{\otimes L})\geq1-(\operatorname{sech}J)^L.
\end{equation}
The spectral constant stays fixed while both global metric coefficients
degenerate.  This is the promised no-tension example: the obstruction belongs
to the global instruments, not to slow mixing of the independent system.

\section{From a link to the Dobrushin window}

Let $\gamma_i(\cdot\mid x_{-i})$ be one-site conditional distributions and
let $C=(c_{ij})$ be a Dobrushin interdependence matrix, so changing only site
$j$ changes the conditional at $i$ in total variation by at most $c_{ij}$.
The classical comparison theorem propagates local oscillations through $C$;
see \cite{Dobrushin1968,Dobrushin1970} and the cone formulation in
\cite{GaubertQu2014}.

If
\begin{equation}\label{eq:alpha-window}
 \alpha=\sup_i\sum_jc_{ij}<1,
\end{equation}
then the Neumann series $D=(I-C)^{-1}=\sum_{n\geq0}C^n$ is uniformly bounded
in the row-sum norm by $(1-\alpha)^{-1}$.  If $c_{ij}=0$ when graph distance
$d(i,j)>1$, then
\begin{equation}\label{eq:spatial-resolvent}
 D_{ij}\leq\sum_{n\geq d(i,j)}(C^n)_{ij}
 \leq\frac{\alpha^{d(i,j)}}{1-\alpha}.
\end{equation}
This is where volume-uniformity lives: in a local vector inequality plus a
subcritical row sum, not in the full tensor-cone diameter.

For a pairwise Ising specification, Proposition~\ref{prop:ising} gives the
uniform envelope
\[
 c_{ij}\leq\tanh|J_{ij}|.
\]
The supremum over the ambient real field is sharp, although the maximizing
field need not be attainable in a particular finite boundary-value problem.
On an anisotropic square lattice with horizontal coupling $\beta$ and vertical
coupling $\gamma$, the familiar sufficient window is
\begin{equation}\label{eq:anisotropic-window}
 2\tanh|\beta|+2\tanh|\gamma|<1.
\end{equation}
It is sufficient and generally not necessary.  Nothing here extends a result
beyond the Dobrushin window or proves optimality of that window.

\section{What the wall rules out---and what it does not}

\begin{theorem}[Comparative theorem with its wall]\label{thm:main-wall}
For every strictly positive finite Markov kernel $P$:
\begin{enumerate}[label=(\arabic*)]
\item Birkhoff's coefficient on the chosen global cone is exactly
$\tau(P)=\tanh(\diamH(P)/4)$.
\item Dobrushin's coefficient is exactly the oscillation contraction factor
and satisfies $\delta(P)\leq\tau(P)$, with equality precisely as in
Theorem~\ref{thm:kernel-equality}.
\item Under independent tensor powers, $\tau(P^{\otimes L})$ follows
\eqref{eq:tensor-power} and tends to one unless $P$ has identical rows.
The global $\delta(P^{\otimes L})$ also tends to one unless all rows agree.
\item Each coordinate oscillation contracts with exact factor $\delta(P)$,
uniformly in $L$.  In interacting systems the analogous uniform comparison
is controlled by \eqref{eq:alpha-window}.
\item For the binary Ising cell the three local constants---projective,
Markov--Dobrushin, and conditional influence envelope---all equal
$\tanh|J|$, but the global projective tensor coefficient equals
$\tanh(L|J|)$.
\end{enumerate}
Therefore no volume-uniform gap follows from applying the global Birkhoff
diameter calculation to the tensor cone.  This obstruction is part of the
comparative conclusion.
\end{theorem}

The theorem rules out the proposed mechanism, not every projective method.
A genuinely local cone, a quotient removing extensive directions, or a
nonproduct Finsler structure might behave differently.  Such a proposal must
specify its cone and metric and pass the tensor test before a formalization
campaign.  Merely inserting positive diagonal Gibbs weights does not count,
by Proposition~\ref{prop:diagonal}.

\section{Reproducibility and formalization status}

The companion script
\texttt{scripts/killtest\_birkhoff\_tensor\_wall.py} uses exact rational
arithmetic for
\[
 K=\begin{pmatrix}3&1\\1&3\end{pmatrix}.
\]
It verifies by enumeration through $L=4$ that
\[
 \Theta(K^{\otimes L})=9^L,\qquad
 \tau(K^{\otimes L})=\frac{3^L-1}{3^L+1},
\]
whereas the normalized single-excitation eigenvalue remains $1/2$.  A second
exact-arithmetic certificate checks the pairwise envelope, equality and
strict examples, cross-ratio tensorization, and coordinatewise saturation.
These computations are regression tests, not substitutes for the proofs.

The companion module
\texttt{YangMills/BirkhoffDobrushin/Wall.lean} formalizes the algebraic core
that survived the kill test.  In particular it checks: the scalar square
identity and reciprocal-extrema equality condition; nonnegativity and endpoint
rigidity of the chord slack, including its strictly positive finite weighted
sum form; factorization of projective cross-ratios and preservation of a
realized maximum under tensor products; the rational hyperbolic-addition law
for Birkhoff coefficients; and the finite conditioning estimate that carries
a fibrewise oscillation bound through averaging.  The exact source was
compiled with Lean \texttt{v4.29.0-rc6} on a Colab Pro+ CPU high-memory
runtime.  The successful package-import artifacts are:
\begin{center}
\scriptsize
\begin{tabular}{@{}ll@{}}
Mathlib commit & \texttt{07642720480157414db592fa85b626dafb71355b} \\
source SHA-256 & \texttt{a477ca8b77133583b16f40f9a55ffab6da6296d873a62de7bebddd1571dca0e8} \\
module SHA-256 & \texttt{1283835e606faa1797d86d47ea01b23e1aa1baf378a9c3d559642f6c60026490}
\end{tabular}
\end{center}
Both module compilation and a separate importing file had exit code zero.

This is deliberately not described as an end-to-end formalization of every
analytic statement in the paper.  Probability normalization, the general
Birkhoff--Hopf contraction theorem, the logarithmic passage from maximal
cross-ratio to Hilbert diameter, the Hellinger proof for the global product
Dobrushin coefficient, and the interacting Dobrushin comparison theorem remain
classical external inputs.  Failed compiler and packaging passes are retained
in the verification transcript rather than erased.  No Lean or Lake process
was run on the Windows owner desktop.

\section{Limitations and conclusion}

The analysis is finite and strictly positive.  Zeros require faces of the
cone or extended projective distances.  The Dobrushin window is not claimed
optimal.  The local product theorem concerns independent coordinate updates;
interacting dynamics require a specified update rule and comparison matrix.
No Yang--Mills mass gap, infinite-volume spectral gap, or result outside the
Dobrushin regime follows from this note.

Within those limits, the conclusion is sharp.  The local binary cell explains
the common $\tanh$.  Equality is completely visible in a two-valued reciprocal
likelihood ratio.  Tensorization then turns the same hyperbolic parameter into
an extensive quantity and drives the global coefficient to one.  Dobrushin's
uniformity survives only after retaining coordinatewise information.  The
calculation and the wall are two sides of the same theorem.

\begin{thebibliography}{99}

\bibitem{Birkhoff1957}
G.~Birkhoff,
\newblock Extensions of Jentzsch's theorem,
\newblock \emph{Transactions of the American Mathematical Society} 85
(1957), 219--227.
\newblock \href{https://doi.org/10.1090/S0002-9947-1957-0087058-6}
{doi:10.1090/S0002-9947-1957-0087058-6}.

\bibitem{BushellRatio}
P.~J. Bushell,
\newblock On the projective contraction ratio for positive linear mappings,
\newblock \emph{Journal of the London Mathematical Society} (2) 6 (1973),
256--258.
\newblock \href{https://doi.org/10.1112/jlms/s2-6.2.256}
{doi:10.1112/jlms/s2-6.2.256}.

\bibitem{BushellMetric}
P.~J. Bushell,
\newblock Hilbert's metric and positive contraction mappings in a Banach
space,
\newblock \emph{Archive for Rational Mechanics and Analysis} 52 (1973),
330--338.
\newblock \href{https://doi.org/10.1007/BF00247467}
{doi:10.1007/BF00247467}.

\bibitem{CohenFausti2024}
S.~Cohen and M.~Fausti,
\newblock Hyperbolic contractivity and the Hilbert metric on probability
measures,
\newblock arXiv:2309.02413v2 (2024).
\newblock \href{https://arxiv.org/abs/2309.02413}{arXiv:2309.02413}.

\bibitem{Dobrushin1956}
R.~L. Dobrushin,
\newblock Central limit theorem for nonstationary Markov chains. I,
\newblock \emph{Theory of Probability and Its Applications} 1 (1956),
65--80.
\newblock \href{https://doi.org/10.1137/1101006}{doi:10.1137/1101006}.

\bibitem{Dobrushin1968}
R.~L. Dobrushin,
\newblock The description of a random field by means of conditional
probabilities and conditions of its regularity,
\newblock \emph{Theory of Probability and Its Applications} 13 (1968),
197--224.
\newblock \href{https://doi.org/10.1137/1113026}{doi:10.1137/1113026}.

\bibitem{Dobrushin1970}
R.~L. Dobrushin,
\newblock Prescribing a system of random variables by conditional
distributions,
\newblock \emph{Theory of Probability and Its Applications} 15 (1970),
458--486.
\newblock \href{https://doi.org/10.1137/1115049}{doi:10.1137/1115049}.

\bibitem{GaubertQu2014}
S.~Gaubert and Z.~Qu,
\newblock Dobrushin's ergodicity coefficient for Markov operators on cones,
\newblock \emph{Integral Equations and Operator Theory} 81 (2015), 127--150.
\newblock \href{https://doi.org/10.1007/s00020-014-2193-2}
{doi:10.1007/s00020-014-2193-2}.

\bibitem{ReebKastoryanoWolf2011}
D.~Reeb, M.~J. Kastoryano and M.~M. Wolf,
\newblock Hilbert's projective metric in quantum information theory,
\newblock \emph{Journal of Mathematical Physics} 52 (2011), 082201.
\newblock \href{https://doi.org/10.1063/1.3615729}
{doi:10.1063/1.3615729}.

\end{thebibliography}

\end{document}
